\chapter*{Термины и определения}
\addcontentsline{toc}{chapter}{Термины и определения}

В настоящем отчёте применяют следующие термины с соответствующими определениями.

Прямой рендеринг -- процесс получения изображения 3D-сцены по заданным моделям геометрии, материалам и источникам освещения.

Обратный рендеринг -- процесс оптимизации (восстановления) неизвестных компонент 3D-сцены (одной или нескольких) по фотографиям и некоторым известным компонентам сцены.

Дифференцируемый рендеринг -- процесс вычисления градиента функции изображения по параметрам 3D-сцены.

Функция расстояния со знаком -- скалярная функция, значение которой в точке пространства равно расстоянию до поверхности объекта, взятому со знаком минус внутри объекта и со знаком плюс снаружи.

Брик -- небольшая регулярная сетка значений функции расстояния со знаком, хранящаяся в листе иерархии ограничивающих объёмов.

Разреженное множество бриков -- представление функции расстояния со знаком, в котором хранятся только брики, содержащие поверхность объекта, и брики в их окрестности.

Редистансинг -- пересчёт значений сетки функции расстояния со знаком, восстанавливающий их соответствие расстояниям до поверхности, заданной нулевым уровнем этой функции.

Релаксированная граница -- множество точек вдоль луча, в которых значение функции расстояния со знаком меньше заданного порога и достигает локального минимума; используется для оценки вклада изменения видимости в градиент.

Эпоха реконструкции -- последовательность итераций оптимизации при фиксированном разрешении сетки функции расстояния со знаком.

Гониометрическая диаграмма направленности -- зависимость силы света (яркости) источника от направления излучения, заданная таблицей значений по сфере направлений.

Кэш световых путей -- сохранённый набор световых путей сцены с их весами, позволяющий вычислять изображение и его градиент по параметрам излучения источника без повторной трассировки лучей.

\chapter*{Обозначения и сокращения}
\addcontentsline{toc}{chapter}{Обозначения и сокращения}

ИИ -- искусственный интеллект;

МНК -- метод наименьших квадратов;

ПО -- программное обеспечение;

ТЗ -- техническое задание;

Adam -- ADAptive Moment estimation, вариант градиентного спуска с адаптивным выбором шага;

BVH -- Bounding Volume Hierarchy, иерархия ограничивающих объёмов;

CPU -- Central Processing Unit, центральный процессор;

GPU -- Graphics Processing Unit, графический процессор;

HDR -- High Dynamic Range, расширенный динамический диапазон;

IES -- формат фотометрических данных светильников (стандарт IES LM-63);

LDR -- Low Dynamic Range, ограниченный динамический диапазон;

MSE -- Mean Squared Error, среднеквадратичная ошибка;

PRB -- Path Replay Backpropagation, обратное распространение с повторным воспроизведением путей;

PSNR -- Peak Signal-to-Noise Ratio, пиковое отношение сигнала к шуму;

SBS -- Sparse Brick Set, разреженное множество бриков;

SDF -- Signed Distance Function, функция расстояния со знаком.

\endinput
